% 脉冲星行星（综述）
% license CCBYSA3
% type Wiki

本文根据 CC-BY-SA 协议转载翻译自维基百科\href{https://en.wikipedia.org/wiki/Pulsar_planet}{相关文章}。

\begin{figure}[ht]
\centering
\includegraphics[width=6cm]{./figures/427823341c42e767.png}
\caption{带有行星的脉冲星艺术概念图} \label{fig_Pupl_1}
\end{figure}
脉冲星行星是指围绕脉冲星运行的行星。1992 年，人们首次在一颗毫秒脉冲星周围发现了这类行星，它们也是人类确认发现的第一批太阳系外行星。脉冲星本身是极其精确的“宇宙时钟”，即使质量很小的行星，也能在脉冲星的观测特性中引起可探测的变化；目前已知质量最小的系外行星，正是一颗脉冲星行星。

这类行星极为罕见，NASA 系外行星档案中仅列出大约半打（六颗左右）。只有一些非常特殊的过程，才能在脉冲星周围形成行星级别的伴星，其中许多被认为是性质十分奇特的天体，例如由伴星部分被破坏后形成的“钻石行星”。此外，脉冲星强烈的辐射以及由电子–正电子对组成的高速粒子风，往往会剥离这些行星的大气层，因此它们几乎不可能成为生命适宜的栖居地。
\subsection{形成}
行星的形成需要原行星盘的存在，而多数理论还要求盘内存在一个“**死区**”，即没有湍流的区域。在该区域内，微行星体可以形成并不断积累，而不会因向恒星内落而被吞噬 $^{[1]}$。与年轻恒星相比，脉冲星具有高得多的光度，因此其辐射会强烈电离盘物质，从而抑制死区的形成 $^{[2]}$；电离会触发**磁旋转不稳定性**，引发湍流并破坏死区 $^{[3]}$。因此，若要孕育行星，围绕脉冲星的盘必须具有相当大的质量 $^{[4]}$。

目前提出了若干可能形成行星系统的过程 $^{[a]}$：

\begin{itemize}
\item “第一代”行星：指在恒星发生超新星爆发并演化为脉冲星之前，就已围绕其运行的行星 $^{[6]}$。大质量恒星往往缺乏行星，这可能既源于在极其明亮的恒星周围难以探测行星，也由于强辐射会破坏原行星盘。若行星轨道半径小于约 4 个天文单位，当恒星演化为红巨星或红超巨星时，行星很可能被吞没并毁灭。在超新星爆发过程中，系统会损失大约一半的质量；除非脉冲星在爆发时被“踢出”的方向恰与行星当时的运动方向一致，否则行星往往会与系统解体而脱离。已知的脉冲星行星系统中，没有一个很可能是通过这一过程形成的 $^{[7]}$。

\item “第二代”行星**：由超新星爆发后回落到脉冲星上的物质形成 $^{[6]}$。理论上，这些回落物质的总质量可与典型原行星盘相当 $^{[7]}$，但其耗散速度可能过快，不利于行星形成。目前尚无围绕年轻脉冲星的已知行星实例 $^{[8][9]}$。

\item “第三代”行星** $^{[6]}$：脉冲星与伴星相互作用并将其摧毁，从而形成一个低质量盘。脉冲星可发出高能辐射加热伴星，使其充满并溢出罗希瓣，最终被完全破坏。另一种机制是引力波辐射导致轨道不断收缩，直到伴星（在这些情形中通常为白矮星）解体 $^{[8]}$。第三种机制是脉冲星进入一颗更大恒星的包层，使其解体并在脉冲星周围形成一个盘 $^{[10][11]}$。由这些过程形成的盘，其质量远大于回落物质形成的盘，因此能够存活更长时间，从而允许行星形成 $^{[8]}$。这些盘还富含行星形成所必需的重元素，并且盘物质的一部分会被脉冲星吸积，从而使其自转加快 $^{[12]}$。此外，也可能出现另一种情形：一颗较轻的白矮星在与更大质量白矮星的相互作用中被摧毁，轻白矮星形成的碎屑盘孕育行星，而较大的白矮星最终演化为脉冲星 $^{[13]}$。
\item 在某些情况下，伴星在与脉冲星的相互作用中被摧毁，但会留下一个行星尺度的残余体 $^{[3]}$。这类系统被称为“黑寡妇”系统$^{[14]}$，但并不总被视为脉冲星行星的一种 $^{[15]}$。
\item 最后，还有可能是来自伴星的行星或流浪行星被脉冲星俘获 $^{[16]}$，或者脉冲星与行星原先所围绕的宿主恒星发生并合 $^{[17]}$。后一种过程会形成一个公共包层，该包层最终瓦解并形成一个盘，从中孕育行星 $^{[18]}$。
\end{itemize}